\documentclass{fiche}
\metadonnees{11}{Coketown}{Bloc 3 — La ville, le travail, la science}

\begin{document}
\entete

\objectifs{%
\textbf{Langue} : déterminants (Ø / \emph{the} / \emph{a}) ; génitif ; noms composés.\par
\textbf{Texte} : Charles Dickens, \emph{Hard Times} (1854), livre \textsc{i}, chapitre
\textsc{v}.\par
\textbf{Durée} : 1\,h\,30 — exercices 35 min, texte et questions 30 min, version 25 min.}

%% ===================================================================
\exercice[13 min]{Ø, \emph{the} ou \emph{a}}

\rappel{\textbf{Rappel.} Pas d'article devant un nom pris en général, qu'il soit pluriel ou
indénombrable. \emph{The} renvoie à ce qui est déjà connu ou déterminé par la suite de la
phrase.}

\consigne{Complétez par \emph{the}, \emph{a} ou rien (Ø).}

\begin{anglais}
\begin{enumerate}[label=\arabic*.,itemsep=2.2mm,leftmargin=*]
  \item \trou{Ø} smoke covered the whole town.
  \item \trou{The} smoke of the tall chimneys never stopped.
  \item It was \trou{a} town of red brick.
  \item \trou{Ø} bricks were red under \trou{the} soot.
  \item \trou{Ø} machinery does not rest.
  \item \trou{The} machinery of that mill is fifty years old.
  \item \trou{Ø} work begins at six and ends at seven.
  \item He found \trou{Ø} work in a warehouse.
  \item \trou{Ø} fact was all that mattered to Mr Gradgrind.
  \item \trou{The} facts he wanted were figures.
  \item \trou{Ø} life in Coketown was the same every day.
  \item \trou{The} life of a factory hand was short.
  \item \trou{Ø} men and \trou{Ø} women went in and out at \trou{the} same hours.
  \item It is \trou{a} folly to argue with \trou{Ø} facts.
\end{enumerate}
\end{anglais}

\note{Sans article : 1, 4, 5, 7, 8, 9, 11, 13 (\emph{men}, \emph{women}). Avec \emph{the} :
2, 6, 10, 12, 13 (\emph{the same hours}).}

\note[Le test]{Si le nom peut être suivi de \og en général \fg, pas d'article :
\emph{smoke covered the town} (la fumée en général). S'il est suivi d'un complément ou d'une
relative qui le délimite, \emph{the} s'impose : \emph{the smoke of the tall chimneys},
\emph{the facts he wanted}. C'est la règle la plus rentable de tout le programme, et celle
qui coûte le plus de points aux francophones, dont la langue met \og le \fg{} partout.}

%% ===================================================================
\exercice[10 min]{Le génitif}

\rappel{\textbf{Rappel.} Le génitif en \emph{'s} s'emploie pour les êtres animés, les
collectivités et les expressions de temps ; \emph{of} pour les choses.}

\consigne{Reformulez avec un génitif ou avec \emph{of}, selon ce qui convient.}

\begin{anglais}
\begin{enumerate}[label=\arabic*.,itemsep=2.5mm,leftmargin=*]
  \item the house of my sister \trou{my sister's house}
  \item the top of the chimney \trou{the top of the chimney}
  \item the work of a day \trou{a day's work}
  \item the face of the savage \trou{the savage's face}
  \item the windows of the mill \trou{the windows of the mill}
  \item the school of Mr Gradgrind \trou{Mr Gradgrind's school}
  \item the noise of the machines \trou{the noise of the machines}
  \item the holidays of the children \trou{the children's holidays}
  \item the journey of an hour \trou{an hour's journey}
  \item the streets of the town \trou{the town's streets \emph{ou} the streets of the town}
\end{enumerate}
\end{anglais}

\note[Trois précisions]{Pluriel en \emph{-s} : l'apostrophe seule (\emph{the workers'
wages}) ; pluriel irrégulier : \emph{'s} (\emph{the children's holidays}). Une ville, un
pays, une institution se comportent comme des animés : \emph{the town's streets},
\emph{France's history} sont corrects à côté de \emph{of}. Enfin le génitif se combine avec
l'article du possesseur, jamais avec celui du possédé : \emph{my sister's house}, et non
\emph{the my sister's house}.}

%% ===================================================================
\exercice[12 min]{Noms composés}

\rappel{\textbf{Rappel.} Dans un nom composé, le dernier mot porte le sens principal ; ce qui
précède le qualifie. On lit donc de droite à gauche.}

\consigne{a) Traduisez.}

\begin{multicols}{2}
\begin{anglais}
\begin{enumerate}[label=\arabic*.,itemsep=1.4mm,leftmargin=*]
  \item a steam-engine \trou{une machine à vapeur}
  \item a factory hand \trou{un ouvrier d'usine}
  \item a brick wall \trou{un mur de briques}
  \item a town hall \trou{un hôtel de ville}
  \item a school-room \trou{une salle de classe}
  \item a chimney sweep \trou{un ramoneur}
\end{enumerate}
\end{anglais}
\columnbreak
\begin{anglais}
\begin{enumerate}[label=\arabic*.,itemsep=1.4mm,leftmargin=*,start=7]
  \item a coal mine \trou{une mine de charbon}
  \item a cotton mill \trou{une filature de coton}
  \item a railway station \trou{une gare}
  \item a warehouse \trou{un entrepôt}
  \item a bird cage \trou{une cage à oiseaux}
  \item a water pipe \trou{une conduite d'eau}
\end{enumerate}
\end{anglais}
\end{multicols}

\note{Le français inverse l'ordre et ajoute une préposition : \emph{a coal mine}
$\rightarrow$ une mine \emph{de} charbon. Traduire de droite à gauche donne toujours le bon
noyau. Le premier élément reste au singulier même quand le sens est pluriel : \emph{a
three-year-old child}, \emph{a ten-mile walk}.}

\consigne{b) Formez le nom composé correspondant.}
\begin{anglais}
\begin{enumerate}[label=\arabic*.,itemsep=2mm,leftmargin=*]
  \item une marche de deux heures \trou{a two-hour walk}
  \item une cheminée d'usine \trou{a factory chimney}
  \item un enfant de dix ans \trou{a ten-year-old child}
  \item un ouvrier du textile \trou{a textile worker}
\end{enumerate}
\end{anglais}

%% ===================================================================
\newpage
\section{Texte}

\noindent\small\textit{Coketown est une ville industrielle imaginaire du nord de
l'Angleterre, bâtie sur le charbon et le coton.}\normalsize

\begin{texte}
It was a town of red brick, or of brick that would have been red if the smoke and
\gl[ashes]{ashes}{les cendres} had allowed it; but as matters stood, it was a town of
unnatural red and black like the painted face of a \gl[a savage]{savage}{un sauvage}. It was
a town of machinery and tall chimneys, out of which interminable serpents of smoke trailed
themselves for ever and ever, and never got \gl[to uncoil]{uncoiled}{se dérouler,
se déployer}. It had a black canal in it, and a river that ran purple with ill-smelling
\gl[dye]{dye}{la teinture}, and vast piles of building full of windows where there was a
rattling and a trembling all day long, and where the piston of the steam-engine worked
monotonously up and down, like the head of an elephant in a state of melancholy madness. It
contained several large streets all very like one another, and many small streets still more
like one another, inhabited by people equally like one another, who all went in and out at
the same hours, with the same sound upon the same pavements, to do the same work, and to whom
every day was the same as yesterday and to-morrow, and every year the
\gl[a counterpart]{counterpart}{le pendant, la réplique} of the last and the next.

These attributes of Coketown were in the main inseparable from the work by which it was
sustained; against them were to be set off, comforts of life which found their way all over
the world, and \gl[elegancies]{elegancies}{les raffinements} of life which made, we will not
ask how much of the fine lady, who could scarcely bear to hear the place mentioned. The rest
of its features were voluntary, and they were these.

You saw nothing in Coketown but what was severely \gl[workful]{workful}{voué au travail ---
mot forgé par Dickens}. If the members of a religious \gl[a persuasion]{persuasion}{une
confession, une obédience} built a chapel there --- as the members of eighteen religious
persuasions had done --- they made it a pious warehouse of red brick, with sometimes (but
this is only in highly ornamental examples) a bell in a birdcage on the top of it. The
solitary exception was the New Church; a \gl[stuccoed]{stuccoed}{crépi, stuqué} edifice with
a square \gl[a steeple]{steeple}{un clocher} over the door, terminating in four short
\gl[a pinnacle]{pinnacles}{un pinacle} like \gl[florid]{florid}{ornementé, chantourné} wooden
legs. All the public inscriptions in the town were painted alike, in severe characters of
black and white. The jail might have been the \gl[an infirmary]{infirmary}{un hospice,
un dispensaire}, the infirmary might have been the jail, the town-hall might have been
either, or both, or anything else, for anything that appeared to the contrary in the graces
of their construction. Fact, fact, fact, everywhere in the material aspect of the town; fact,
fact, fact, everywhere in the immaterial. The M'Choakumchild school was all fact, and the
school of design was all fact, and the relations between master and man were all fact, and
everything was fact between the \gl[a lying-in hospital]{lying-in hospital}{une maternité}
and the cemetery, and what you couldn't state in figures, or show to be
\gl[purchaseable]{purchaseable}{achetable} in the cheapest market and
\gl[saleable]{saleable}{vendable} in the dearest, was not, and never should be, world without
end, Amen.
\end{texte}
\source{Charles Dickens, \emph{Hard Times}, Londres, Bradbury \& Evans, 1854, livre
\textsc{i}, ch.~\textsc{v}.}

\partiel[15 min]{Questions}
\consigne{\en{Answer in English, in full sentences.}}

\begin{enumerate}[label=\arabic*.,itemsep=2mm,leftmargin=*]
  \item \en{What colours is Coketown made of, and why?}
    \lignes{2}
    \corr{Red and black: the brick would have been red, but smoke and ashes have blackened
    it, so the town looks like the painted face of a savage.}
  \item \en{Find the three comparisons in the first paragraph and say what each one adds.}
    \lignes{3}
    \corr{The smoke as interminable serpents that never uncoil (something endless and
    threatening); the town as a painted savage (a civilisation that has turned barbarous);
    the piston as an elephant's head in melancholy madness (a huge patient animal driven mad
    by repetition).}
  \item \en{What is repeated in the last sentence of the first paragraph, and what is the
    effect?}
    \lignes{3}
    \corr{\emph{The same} and \emph{like one another} come back again and again: the same
    hours, sound, pavements, work, every day like yesterday and to-morrow. The sentence
    itself becomes as monotonous as the life it describes.}
  \item \en{How can you tell the chapels, the jail, the infirmary and the town hall apart?}
    \lignes{2}
    \corr{You cannot. They were all built alike, and nothing in their construction says what
    they are for.}
  \item \en{What does the town consider real?}
    \lignes{3}
    \corr{Only what can be stated in figures or bought cheap and sold dear. Anything that
    cannot be counted or traded \emph{was not, and never should be}.}
  \item \en{The paragraph ends with ``world without end, Amen.'' Why does Dickens end on
    the words of a prayer?}
    \lignes{3}
    \corr{Because fact has taken the place of religion in Coketown. The chapels are
    warehouses, and the real creed is the one he has just stated; ending it like a prayer
    makes the irony complete.}
\end{enumerate}

%% ===================================================================
\newpage
\section{Version}
\consigne{Traduisez le passage en français. C'est l'ouverture du roman : un notable expose
sa doctrine au maître d'école.}

\begin{version}
`Now, what I want is, Facts. Teach these boys and girls nothing but Facts. Facts alone are
wanted in life. Plant nothing else, and \gl[to root out]{root out}{extirper, déraciner}
everything else. You can only form the minds of reasoning animals upon Facts: nothing else
will ever be of any service to them. This is the principle on which I bring up my own
children, and this is the principle on which I bring up these children. Stick to Facts, sir!'

The scene was a plain, bare, monotonous \gl[a vault]{vault}{une voûte, une cave} of a
school-room, and the speaker's square \gl[a forefinger]{forefinger}{l'index} emphasized his
observations by \gl[to underscore]{underscoring}{souligner} every sentence with a line on the
schoolmaster's sleeve. The emphasis was helped by the speaker's square wall of a forehead,
which had his eyebrows for its base, while his eyes found
\gl[commodious]{commodious}{spacieux} \gl[cellarage]{cellarage}{un espace de cave} in two dark
caves, overshadowed by the wall. The emphasis was helped by the speaker's mouth, which was
wide, thin, and hard set. The emphasis was helped by the speaker's voice, which was
inflexible, dry, and dictatorial. The emphasis was helped by the speaker's hair, which
\gl[to bristle]{bristled}{se hérisser} on the skirts of his bald head, a plantation of
\gl[a fir]{firs}{un sapin} to keep the wind from its shining surface, all covered with
\gl[a knob]{knobs}{une bosse, un nœud}, like the crust of a plum pie, as if the head had
scarcely warehouse-room for the hard facts stored inside.
\end{version}
\source{\emph{Ibid.}, livre \textsc{i}, ch.~\textsc{i}.}

\lignes{22}

\corr{\textbf{Proposition de traduction.} \og Voilà ce que je veux : des Faits. N'enseignez à
ces garçons et à ces filles rien que des Faits. Seuls les Faits sont requis dans la vie. Ne
plantez rien d'autre, et arrachez tout le reste. On ne peut former l'esprit d'animaux
raisonnables que sur des Faits : rien d'autre ne leur sera jamais de la moindre utilité.
C'est le principe sur lequel j'élève mes propres enfants, et c'est le principe sur lequel
j'élève ces enfants-ci. Tenez-vous-en aux Faits, monsieur ! \fg

La scène se passait dans une salle de classe nue, unie, monotone comme une cave, et l'index
carré de l'orateur appuyait ses observations en soulignant chaque phrase d'un trait sur la
manche du maître d'école. L'insistance était secondée par le mur carré qu'était le front de
l'orateur, lequel avait ses sourcils pour soubassement, tandis que ses yeux trouvaient de
spacieux caveaux dans deux grottes sombres que ce mur surplombait. L'insistance était
secondée par la bouche de l'orateur, large, mince et serrée. L'insistance était secondée par
la voix de l'orateur, inflexible, sèche et dictatoriale. L'insistance était secondée par les
cheveux de l'orateur, qui se hérissaient sur les bords de son crâne chauve, plantation de
sapins destinée à protéger du vent sa surface luisante, toute couverte de bosses comme la
croûte d'une tourte aux prunes, comme si la tête n'avait pas assez de place d'entrepôt pour
les faits bien durs qu'elle emmagasinait.}

\note[Choix de traduction 1]{\emph{a vault of a school-room}, \emph{his square wall of a
forehead} : tour en \emph{N of a N} où le premier nom est en réalité une métaphore appliquée
au second. Le français n'a pas cette construction ; il passe par \og comme \fg{} ou par une
relative (\og le mur carré qu'était son front \fg).}

\note[Choix de traduction 2]{\emph{Facts alone are wanted in life} : \emph{to want} au passif
a le sens ancien de \og être requis, faire défaut \fg. Le contexte l'impose : ce n'est pas
\og les faits sont voulus \fg.}

\note[Choix de traduction 3]{\emph{The emphasis was helped by\ldots} répété quatre fois : le
passif sert ici à garder le même sujet en tête de phrase et à faire défiler les
compléments. Le français peut garder le passif, à condition de varier le verbe le moins
possible : c'est l'effet de litanie qui compte.}

\note[Choix de traduction 4]{\emph{Stick to Facts} : \emph{to stick to} = s'en tenir à.
Attention à l'impératif français : \og tenez-vous-en \fg, avec les deux traits d'union.}

%% ===================================================================
\partiel{Lexique à mémoriser}
\begin{itemize}[label=--,itemsep=0.8mm,leftmargin=*]\small
  \item \textbf{a brick} — une brique ; a brick wall
  \item \textbf{smoke} — la fumée \textit{(indénombrable)} ; \textit{verbe} to smoke
  \item \textbf{ashes} — les cendres ; \textit{singulier} ash
  \item \textbf{soot} — la suie
  \item \textbf{a chimney} — une cheminée (le conduit) ; \textit{à distinguer de} a fireplace \textit{(l'âtre)}
  \item \textbf{machinery} — les machines \textit{(indénombrable collectif)}
  \item \textbf{a mill} — une usine textile, un moulin ; a cotton mill
  \item \textbf{a warehouse} — un entrepôt ; ware \textit{« la marchandise » +} house
  \item \textbf{a hand} — un ouvrier ; \textit{de} factory hand\textit{, la main qui travaille}
  \item \textbf{a pavement} — un trottoir ; \textit{en anglais américain} a sidewalk
  \item \textbf{a jail} — une prison ; \textit{variante} gaol\textit{, même prononciation}
  \item \textbf{a steeple} — un clocher
  \item \textbf{alike} — semblable, de la même façon ; a- + like
  \item \textbf{to sustain} — soutenir, faire vivre
  \item \textbf{to bear} — supporter ; she could scarcely bear to hear it
  \item \textbf{scarcely} — à peine ; \textit{proche de} hardly
  \item \textbf{a feature} — un trait, une caractéristique
  \item \textbf{figures} — les chiffres ; \textit{et non « les figures »} --- \textit{faux ami}
  \item \textbf{cheap} — bon marché ; \textit{contraire} dear
  \item \textbf{to state} — énoncer, déclarer ; \textit{nom} a statement
\end{itemize}

\end{document}
